\section{Impedance velocity control}
\subsection{Force/torque sensor filtering}
Combination of moving window average and low-pass filter.
TODO: put in the equation

\paragraph{Wrench frame change}
TODO: write in more detail and

Use
\begin{equation}
\begin{bmatrix} 
		R & 0 \\
		0 & R
\end{bmatrix} 
\end{equation}
because you only want frame change, not scaling!
However, if you want the compute a \textit{different wrench},
the one experience at a different part of the same rigid body
where from which you received the reading,
and still expressed in the same frame,
use the $ Ad T $ operation.

\subsection{Direct force feedback}
As simple as adding $ K_{ f }J^{ T } Z \mathcal{F}  $ to Cartesian objective,
where:
\begin{itemize}
		\item $ K_{ f }  $ is the compliance gain
		\item $ \mathcal{F}  $ is the filtered wrench
		\item $ J  $ is the same Jacobian as in the Cartesian objective task 
(the end-effector Jacobian in 99\% of circumstances)
\item $ Z = \text{diag} (z), Z \in \mathbb{R}^{6}  $ is the wrench direction scaling
\end{itemize} 
This then leads to the compliant Cartesian task reference twist:
\begin{equation}
\mathcal{V}_{ \text{compliant} } =  K_{ c }\mathcal{V}_{ \text{error} } + 
K_{ f } J^{ T } \mathcal{F}
\end{equation}

\paragraph{Space to body frame mapping}

\subsection{TODO: Impedance control}
A 6D spring system would have the following dynamics
\begin{equation}
TODO
\end{equation}
From which the impedance reference error can be computed as
\begin{equation}
TODO
\end{equation}
Leading to the following reference velocity:
\begin{equation}
TODO
\end{equation}
